% --- jeśli nie masz w preambule, dodaj:
% \usepackage{booktabs}
% \usepackage{longtable}
% \usepackage{siunitx} % ładne jednostki, np. MiB, %; opcjonalnie
% \sisetup{detect-weight=true, detect-family=true}

\chapter{Eksperyment: analiza ruchu aplikacji Instagram z użyciem PCAPdroid}
\label{chap:eksperyment-instagram}

\section{Cel i założenia}
Celem eksperymentu jest empiryczna charakterystyka komunikacji sieciowej popularnej aplikacji mobilnej (Instagram) pod kątem:
(i) wykorzystywanych protokołów i transportu (HTTP/1.1, HTTP/2, HTTP/3/QUIC, TLS) \cite{RFC8446,RFC9114,RFC9000},
(ii) wzorców zachowania (impulsy podczas interakcji vs.\ aktywność w tle),
(iii) rozproszenia docelowych usług (domen/IP/portów) oraz
(iv) ograniczeń obserwowalności wynikających z QUIC/HTTP/3 i ewentualnego \emph{certificate pinning}.
Rejestracji dokonano narzędziem PCAPdroid (no-root, tryb PCAPNG z kluczami) i poddano obróbce wg metodyki z Rozdz.\ \ref{chap:metodyka}; do inspekcji \texttt{pcapng} wykorzystywano również Wireshark \cite{WiresharkDocs}.

\section{Środowisko i narzędzia}
\begin{itemize}
  \item Urządzenie z Androidem, konto testowe; sieć laboratoryjna (AP na Linuxie).
  \item \textbf{PCAPdroid}: lokalny VPN (bez roota), \emph{TLS Decryption} aktywne, zapis do \texttt{.pcapng} (z kluczami sesji).
  \item Analiza \texttt{pcapng}: Wireshark/tshark (ALPN, SNI, QUIC/UDP~443, identyfikacja HTTP/1.1/2/3) \cite{WiresharkDocs,RFC9114,RFC9000}.
\end{itemize}

\section{Procedura (pojedyncza iteracja)}
Zgodnie z metodyką, jedna iteracja ($\approx$ 15 min) składa się z: 
\begin{enumerate}
  \item \textbf{Baseline OS (2 min)} – app nieuruchomiona; punkt odniesienia dla tła.
  \item \textbf{Cold start (1 min)} – uruchomienie i ustabilizowanie ekranu głównego.
  \item \textbf{Nawigacja (2 min)} – główne ekrany (feed, profil, wyszukiwarka).
  \item \textbf{Media/operacja (1 min)} – krótkie odtwarzanie/przewijanie, wymuszenie większych transferów.
  \item \textbf{Offline/online (4 min)} – 2 min offline + 2 min po reconnect; obserwacja retry/backoff.
  \item \textbf{Tło (10 min)} –  obserwacja po zamknięciu aplikacji.
\end{enumerate}
Motywacja czasów i uzasadnienia znajdują się w Rozdz.\ \ref{chap:metodyka}. QUIC/HTTP/3 (UDP) z natury ogranicza klasyczny MITM; analiza treści HTTP możliwa jest przede wszystkim w wariancie porównawczym z wymuszeniem \emph{fallbacku} do HTTP/2/TCP (blokada UDP/443) \cite{RFC9114}.

\section{Zebrane dane (ta sesja)}
Eksperyment wykonano w dniu 2025-08-18 w oknie \textbf{17{:}56{:}33–18{:}10{:}14} ($\approx$ \SI{13.7}{\minute}). Zarejestrowano \textbf{302} przepływy (\emph{flows}) o łącznym wolumenie \textbf{$\approx$\,28.0 MiB}; \textbf{RX $\approx$\,27.2 MiB}, \textbf{TX $\approx$\,0.85 MiB}.

\paragraph{Udział protokołów (liczba flow):}
QUIC \textbf{61.9\%}, DNS 16.2\%, TLS 11.9\%, TCP 4.0\%, HTTP 2.6\%, HTTPS 1.7\%, UDP 1.7\%.
Dominującym portem docelowym był \textbf{443} (HTTPS/QUIC; 233 flow), dalej \textbf{53} (DNS; 49 flow), \textbf{5222} (TCP; 12 flow) i \textbf{80} (HTTP; 8 flow).

\section{Wyniki ilościowe}
\subsection{Wolumen i liczność wg protokołu}
\begin{table}[htbp]
  \centering
  \caption{Podsumowanie wg protokołu (Instagram, ta sesja)}
  \label{tab:insta-proto}
  \begin{tabular}{lrrrr}
    \toprule
    \textbf{Protokół} & \textbf{Flows} & \textbf{MiB (total)} & \textbf{B wysłane} & \textbf{B odebrane} \\
    \midrule
    QUIC  & 187 & 27.99 & 855{,}220 & 28{,}500{,}436 \\
    HTTPS & 5   & 0.02  & 12{,}180  & 9{,}063 \\
    TCP   & 12  & 0.02  & 15{,}384  & 3{,}699 \\
    DNS   & 49  & 0.00  & 1{,}475   & 2{,}441 \\
    TLS   & 36  & 0.00  & 2{,}160   & 1{,}440 \\
    UDP   & 5   & 0.00  & 837       & 346 \\
    HTTP  & 8   & 0.00  & 480       & 320 \\
    \midrule
    \textbf{Razem} & \textbf{302} & \textbf{28.04} & \textbf{888{,}736} & \textbf{28{,}517{,}745} \\
    \bottomrule
  \end{tabular}
\end{table}

\subsection{Największe cele (IP/port) wg wolumenu}
\begin{table}[htbp]
  \centering
  \caption{Największe kierunki ruchu (wolumen; Instagram, ta sesja)}
  \label{tab:insta-topdst}
  \begin{tabular}{lllrr}
    \toprule
    \textbf{IP docelowe} & \textbf{Port} & \textbf{Proto} & \textbf{Flows} & \textbf{MiB (total)} \\
    \midrule
    185.60.217.63  & 443  & QUIC  & 137 & 27.52 \\
    157.240.251.63 & 443  & QUIC  & 1   & 0.33 \\
    185.60.217.3   & 443  & QUIC  & 35  & 0.08 \\
    185.60.217.20  & 443  & QUIC  & 2   & 0.02 \\
    185.60.217.175 & 5222 & TCP   & 11  & 0.02 \\
    157.240.251.11 & 443  & QUIC  & 1   & 0.02 \\
    185.60.217.28  & 443  & QUIC  & 1   & 0.01 \\
    157.240.0.63   & 443  & QUIC  & 1   & 0.01 \\
    185.60.217.13  & 443  & HTTPS & 2   & 0.01 \\
    157.240.253.17 & 443  & HTTPS & 1   & 0.01 \\
    \bottomrule
  \end{tabular}
\end{table}

\section{Obserwacje jakościowe}
\begin{itemize}
  \item \textbf{Dominacja HTTP/3/QUIC.} Zdecydowana większość ruchu (liczebnie i wolumenowo) to QUIC na 443/UDP, co jest zgodne z trendem migracji warstwy mediów do HTTP/3 \cite{RFC9114,RFC9000}.
  \item \textbf{Kanał utrzymania sesji.} Występuje równoległy kanał na \textbf{TCP/5222} (11--12 flows, mały wolumen), prawdopodobnie do komunikacji w czasie rzeczywistym/utrzymania połączenia.
  \item \textbf{Intensyfikacja przepływów po interakcji, następnie spadek aktywności.} Zliczenia per minuta (niezamieszczone tu dla zwięzłości) pokazują silny pik około 17{:}59 (setki nowych/kończących się flow), po którym następują krótkie sekwencje aktywności w tle ($\approx$ 18{:}05–18{:}06).
  \item \textbf{Niewielki ślad HTTP/1.1/HTTPS.} Jedynie kilka żądań; główna komunikacja i zasoby idą ścieżką QUIC.
\end{itemize}

\subsection{Subiektywne spostrzeżenia}
\paragraph{Obserwowalność treści.}
Ze względu na transport QUIC/UDP, klasyczny MITM dla TLS (po TCP) nie ma zastosowania; nawet z certyfikatem użytkownika wgląd w treści HTTP jest ograniczony \cite{RFC9114}. Analiza opiera się na metadanych (ALPN, SNI, strumienie, wolumeny, czasy) i korelacji zachowania.
\paragraph{Pinning.}
Nawet przy braku aktywnie zaimplementowanego pinningu, wybór QUIC ogranicza możliwość przechwycenia semantyki HTTP. W aplikacjach z rygorystycznym \emph{certificate pinning} obserwowalność jest dodatkowo zawężona do metadanych (zob.\ \cite{OWASPMASTGPinning}).
\paragraph{Replikowalność i zmienność.}
Wyniki dotyczą konkretnych wersji aplikacji i backendu, daty 2025-08-18, okna ~13.7 min i scenariusza przewijania feedu.

\section{Wariant porównawczy (HTTP/3 $\rightarrow$ HTTP/2)}
Aby uzyskać wgląd w semantykę HTTP, zaplanowano wariant kontrolny z wymuszeniem \emph{fallbacku} z QUIC na HTTP/2/TCP (blokada UDP/443 w sieci labowej) \cite{RFC9114}. Porównywane będą: liczba żądań, rozkład kodów i metod HTTP (gdy odszyfrowalne), TTFB/czasy transferu, liczba i rola domen docelowych.

\section{Artefakty i replikowalność}
Zachowano \texttt{.pcapng} (plik zawiera zarówno pakiety, jak i klucze sesji), dziennik przebiegu (czas kroków), listę wersji aplikacji/OS oraz wyciągi agregatów (tabele \ref{tab:insta-proto} i \ref{tab:insta-topdst}). Schemat nazewnictwa plików: \texttt{YYYYMMDD\_HHMM\_Instagram\_runN\_net=\{default|udp443\_blocked\}\_vApp\_vAndroid.pcapng}.

\section{Podsumowanie}
Eksperyment pokazuje typowy dla nowoczesnych aplikacji wzorzec: dominację HTTP/3/QUIC dla treści (media), lekki kanał sterujący po TCP/5222 oraz okresowe aktywności w tle. W kontekście prywatności i bezpieczeństwa oznacza to, że analiza semantyczna wymaga wariantów kontrolnych (H3$\rightarrow$H2) lub metod instrumentacyjnych; jednocześnie metadane ruchu pozwalają na wiarygodną charakterystykę wzorców komunikacyjnych i rozproszenia dostawców.
